\chapter{Die Metadaten}

\section{Begriffsbestimmung und Klassifikation}

\subsection{Inhaltsdaten versus Metadaten}
Um die Problematik von Metadaten zu verstehen, bedarf es zunächst einer Definition. Allgemein werden Metadaten oft als \glqq Daten über Daten\grqq\  bezeichnet. Das bedeutet, dass es sich dabei um jene Informationen handelt, die ein System benötigt, um die eigentlichen Nutzdaten zu übermitteln oder zu verarbeiten, beziehungsweise um jene, die bei deren Verarbeitung entstehen.
Der Sicherheitsexperte Bruce Schneier behandelt den Unterschied zwischen Daten und Metadaten in seinem Werk \emph{Data und Goliath - Die Schlacht um die Kontrolle unserer Welt} unter anderem mit dem Vergleich zu einem Telefonat: Daten sind der Gesprächsinhalt, während Metadaten die Telefonnummern der beiden Gesprächsteilnehmer, das Datum und die Uhrzeit, und die Dauer des Telefonats beinhalten~\cite[Kap.~2]{schneier2015data}.
Dieser Vergleich deckt sich auch mit der Schweizer Gesetzgebung. Im \emph{Bundesgesetz betreffend die Überwachung des Post- und Fernmeldeverkehrs (BÜPF)} werden Metadaten als sogenannte \glqq Randdaten\grqq\ bezeichnet. Gemäss Art.~8 BÜPF sind dies: die Daten, aus denen hervorgeht, mit wem, wann, wie lange und von wo aus die überwachte Person Verbindung hat oder gehabt hat, sowie technische Merkmale der entsprechenden Verbindung~\cite{buepf_ch}.
Die technische Definition stammt von Robert W. Shirey, welcher im Rahmen seiner Tätigkeit für die \emph{Internet Engineering Task Force} folgende Definition gibt: Metadaten sind beschreibende Informationen über ein Datenobjekt, das heisst Daten über Daten oder Daten-Labels, die andere Daten beschreiben~\cite{rfc4949}.
\subsection{Dimensionen von Metadaten}

\section{Entstehung von Metadaten im Netzwerk}

\subsection{Metadaten auf Protokollebene (Transport \& IP)}

\subsection{Metadaten auf Anwendungsebene}

\section{Das Sicherheitsproblem: Die Aussagekraft von Metadaten}

\subsection{Traffic Analysis und Datenkorrelation}

\subsection{Profilbildung und Beziehungsnetzwerke}

\subsection{Rechtlicher und gesellschaftlicher Kontext}

\section{Konzepte zur Metadaten-Minimierung}

\subsection{Bekannte Lösungsansätze}

\subsection{Das Trilemma: Sicherheit, Effizienz und Usability}

\section{Zusammenfassung}